\documentclass[11pt,a4paper]{ctexart}
\usepackage[margin=2.2cm]{geometry}
\usepackage{booktabs}
\usepackage{array}
\usepackage{listings}
\usepackage{xcolor}
\usepackage{tikz}
\usetikzlibrary{arrows.meta,positioning,shapes.geometric}
\usepackage{fancyhdr}
\usepackage[hidelinks]{hyperref}
\usepackage{titlesec}
\usepackage{enumitem}
\usepackage{tabularx}
\usepackage{float}
\setlist{itemsep=2pt,topsep=3pt,parsep=0pt}

\definecolor{codebg}{gray}{0.96}
\lstset{
  basicstyle=\ttfamily\footnotesize,
  backgroundcolor=\color{codebg},
  breaklines=true,
  columns=fullflexible,
  keepspaces=true,
  frame=single,
  framerule=0.3pt,
  rulecolor=\color{gray!50},
  xleftmargin=4pt,
  aboveskip=6pt,
  belowskip=6pt,
}

\pagestyle{fancy}
\fancyhf{}
\fancyhead[L]{\small 计算概论大作业 · QQCompanion}
\fancyhead[R]{\small\thepage}
\renewcommand{\headrulewidth}{0.3pt}

\titleformat{\section}{\Large\bfseries}{\thesection}{0.6em}{}
\titleformat{\subsection}{\large\bfseries}{\thesubsection}{0.5em}{}

\title{\bfseries QQCompanion：把大模型接进 QQ 的\\长在线智能体}
\author{姓名：Wang Jiazhen \quad 院系：\underline{\hspace{4cm}} \\[4pt]
\small 仓库：\texttt{https://github.com/illenial-eml/QQCompanion} \quad 提交版本：\texttt{v1.0.0}（commit 见仓库）}
\date{2026 年 11 月}

\begin{document}
\maketitle
\thispagestyle{empty}

\begin{abstract}
\noindent
本项目解决一个很具体的问题：\textbf{我想在自己的 QQ 里有一个真正能长期陪着聊天的 AI}。现成的路都走不通——官方机器人平台个人主体无法被搜索和加好友，逆向协议封号风险高，SaaS 机器人没有长期记忆也不能自定义行为。我用 Node.js 从零实现了一个中间层（48 个模块、约 1.27 万行、\textbf{零第三方运行时依赖}），把 OneBot v11 协议端与大模型接起来，并自己解决了「记忆怎么存、群里什么时候该说话、回复怎么不像机器、进程掉了怎么自己爬起来」这些工程问题。课程知识在这里不是装饰：哈希表决定了去重能不能做到 $O(1)$，二分查找决定了提醒调度的轮询开销，环形缓冲与 LRU 决定了程序能不能连续跑几天而不爆内存。项目配有 56 个测试文件、约 1974 条断言，全部通过。本文如实记录设计取舍、性能实测，以及三个「AI 判断错了、靠实验纠回来」的案例。

\vspace{4pt}
\noindent\textbf{关键词：}OneBot v11；大语言模型；哈希表；环形缓冲；BM25；租约锁；指数退避；自然语言日程
\end{abstract}

\tableofcontents
\newpage

%=========================================================
\section{项目背景与目标}
%=========================================================

\subsection{为谁解决什么问题}

使用者只有一个人：我自己。场景也很朴素——我每天挂着 QQ，希望在私聊里有一个「随时接话、记得住事、能帮我记事」的伙伴；在群里它能像正常人一样，该接话时接一句，插不上话就安静。

现有方案的问题很明确：

\begin{table}[H]
\centering\small
\begin{tabularx}{\textwidth}{@{}p{3.2cm}X@{}}
\toprule
方案 & 不足 \\
\midrule
QQ 官方机器人平台 & 个人主体\textbf{不能被搜索、不能被加好友}，只能靠扫码临时会话；被动回复条数受限；功能受平台审核约束 \\
直接逆向 QQ 协议 & 封号风险高，协议随版本变化，维护成本不可控 \\
现成 AI 聊天 SaaS & 无长期记忆、数据在云端、行为不可定制、不能做日程/课表这类本地能力 \\
\bottomrule
\end{tabularx}
\end{table}

\subsection{方案：协议端当「手」，大模型当「脑」，中间层自己写}

我采用 \textbf{OneBot v11 协议 + NapCat 协议端}：协议端负责与 QQ 服务器通信（它本身是成熟的第三方客户端外壳），我这边的程序只跟它说标准 OneBot 事件与 API。这样既不用碰 QQ 私有协议，又保留完全的控制权。

\begin{figure}[h]
\centering
\resizebox{\textwidth}{!}{%
\begin{tikzpicture}[
  node distance=7mm and 11mm,
  box/.style={draw, rounded corners=2pt, fill=blue!6, minimum width=3.4cm, minimum height=1.1cm, align=center, font=\small},
  core/.style={draw, rounded corners=2pt, fill=green!9, minimum width=4.2cm, minimum height=2.2cm, align=center, font=\small},
  arr/.style={-{Latex[length=2mm]}, thick}
]
\node[box] (qq) {QQ 客户端\\（手机 / 电脑）};
\node[box, right=of qq] (nc) {NapCat 协议端\\OneBot v11 / WebSocket};
\node[core, right=of nc] (core) {\textbf{QQCompanion}\\（本项目，Node.js）\\[2pt]\footnotesize 会话调度 · 记忆 · 模型层\\\footnotesize 能力 · 运维};
\node[box, right=of core, fill=red!6] (llm) {大模型 API\\chat / vision};
\node[box, below=9mm of core, fill=gray!10, minimum width=4.2cm] (store) {本地数据 \texttt{data/*.json}\\（不入库，删即清空）};

\draw[arr] (qq) -- node[above,font=\scriptsize]{消息} (nc);
\draw[arr] (nc) -- node[above,font=\scriptsize]{事件} (core);
\draw[arr] (core) -- node[above,font=\scriptsize]{调用} (llm);
\draw[arr, dashed] (core) -- (store);
\end{tikzpicture}%
}
\caption{系统总体架构：一条消息从 QQ 到模型再回到 QQ，全程在本机完成，只有模型调用出网。}
\end{figure}

\subsection{验收指标（写代码之前就定好）}

\begin{enumerate}[leftmargin=1.6em]
\item 私聊发一句话，\textbf{1.5 秒内}出现第一条回复；
\item 长回复要像人一样分成几条「气泡」，且\textbf{不丢字、不乱断}；
\item 群里\textbf{只在被叫到或话题相关时}说话，不刷屏；
\item 记得前几天的对话，且有「长期记忆」可查询；
\item 能把「明天中午 11:40 提醒我开会」变成真提醒；
\item \textbf{连续运行数天}，进程挂了能自己起来，且同一条消息不会回两遍。
\end{enumerate}

上述 6 项在最终版本全部达成（第 6 项的掉线问题另见第 7 节第 1 条）。

%=========================================================
\section{系统设计}
%=========================================================

\subsection{分层与模块边界}

设计上刻意把「协议」「调度」「智能」「数据」「运维」分开，核心原则是：\textbf{任何一层坏掉，不要带着别的层一起死}。

\begin{table}[H]
\centering\small
\begin{tabularx}{\textwidth}{@{}llX@{}}
\toprule
层 & 模块（文件） & 职责 \\
\midrule
接入层 & \texttt{onebot.js}、\texttt{qq.js} & 连接/重连、事件归一化成统一消息结构、发送与限流 \\
调度层 & \texttt{responder.js}（1.37k 行） & 该不该回、回几段、冷却与配额、命令识别 \\
智能层 & \texttt{brain.js}、\texttt{knowledge.js}、\texttt{style.js} & 提示词组装、模型调用、BM25 检索、风格样本 \\
数据层 & \texttt{memory.js}、\texttt{longterm.js}、\texttt{people.js}、\texttt{names.js} & 短期记忆、长期记忆、印象库、备注簿 \\
能力层 & \texttt{reminders.js}、\texttt{timetable.js}、\texttt{weather.js}、\texttt{search.js}、\texttt{vision.js} & 日程、课表、天气、联网、识图 \\
运维层 & \texttt{instance-lock.js}、\texttt{outbox.js}、\texttt{main.js}、\texttt{log.js} & 单实例租约、失败重发、启动与优雅退出、日志 \\
\bottomrule
\end{tabularx}
\end{table}

\subsection{一条消息的完整生命周期}

\begin{enumerate}[leftmargin=1.6em]
\item \textbf{接入}：WebSocket 收到事件 → 归一化（群 / 私聊、发送者、是否 @、图片、被引用消息）→ 用 \texttt{msg\_id} 去重。
\item \textbf{分流}：先跑「命令层」（日程查询、潜水开关、交接模式、印象查询…），命中就直接回，不花模型的钱。
\item \textbf{调度判定}：群聊走关键词 + 冷却 + 配额 + \textbf{消息是否已过期}四道判定；私聊直接回。
\item \textbf{组装上下文}：最近 40 轮历史（每条截 240 字）→ 摘要 → 按需检索知识库 → 关系/印象 → 人设。
\item \textbf{调模型}：群聊用更小 \texttt{maxTokens}、更低温度；私聊可长一点。
\item \textbf{后处理}：去掉 Markdown、补标点、按语义切成多条气泡、可选配一张表情。
\item \textbf{发送}：失败进发件箱，网络恢复后补发（幂等，不重复）。
\end{enumerate}

\subsection{三个自己最满意的机制}

\textbf{(a) 单实例租约锁。} 实测中曾出现「同一条消息回了两遍」，根因是两个进程都认为自己该回。现在进程启动时写 \texttt{data/bot-instance.lock}（含 pid 与时间戳）并每 5 秒更新，发现锁在 20 秒内被别的 pid 更新过就退出。这就是分布式系统里的 lease，只是单机用文件实现。

\textbf{(b) 发件箱（outbox）。} 网络抖动时发送会失败。失败的消息落盘进 \texttt{data/outbox.json}，重连后按顺序补发；补发前用消息指纹去重，保证不重复。这是「至少一次」加幂等，换来「最终不丢」。

\textbf{(c) 过期消息跳过。} 群消息是串行处理的，等轮到某条消息时，话题可能已经过去十几秒。加了 \texttt{staleSkipMs: 7000}：排队超过 7 秒的不再回。这一条把「答非所问」的投诉直接消掉。

%=========================================================
\section{课程知识的具体应用}
%=========================================================

\subsection{数据结构选型表}

\begin{table}[H]
\centering\footnotesize
\begin{tabularx}{\textwidth}{@{}p{2.6cm}p{5.2cm}X@{}}
\toprule
课程内容 & 项目中的用途 & 为什么不是更朴素的写法 \\
\midrule
哈希表 \texttt{dict/set} & 用户→记忆、QQ 号→印象、备注→QQ 号、\texttt{msg\_id} 去重集合 & 去重若用数组 \texttt{includes} 是 $O(n)$，每条消息都要扫全表；\texttt{Set} 平均 $O(1)$ \\
环形缓冲 & 最近 N 轮对话、去重窗口 500 条、群聊热度窗口 & 固定容量、$O(1)$ 摊还，长时间运行内存有界 \\
LRU 缓存 & 识图结果、搜索结果 & 用 \texttt{Map} 的插入有序性实现，命中 $O(1)$，超容量摘最旧；无需第三方库 \\
二分查找 & 在按时间排序的日程/课表里找「下一条该提醒什么」 & 每秒轮询全表是 $O(n)$；排序 + 二分是 $O(\log n)$ 定位 \\
前缀和 / 滑动窗口 & 近 60 秒群热度、每小时配额统计 & 按分钟分桶计数，更新 $O(1)$、求和 $O(1)$（窗口长度固定） \\
堆 vs 排序数组 & 提醒队列 & 规模 $n<100$，排序数组 + 二分在实际常数上优于堆，代码也短——\textbf{选了简单但够用的} \\
贪心 & 长回复切成多条气泡 & 代价函数单调，贪心即最优；上 DP 是过度设计 \\
双指针 & 切分时先找句末标点、再退回弱标点 & 一次线性扫描，不产生临时数组 \\
BM25 / 倒排索引 & 本地知识库问答 & 只对查询命中词算分，$O(\sum_t df_t)$，不扫全库 \\
指数退避 & 重连、限流重试 & 等待 $b\cdot r^k$（$r=1.4/1.8$，上限 60\,s），避免重试风暴 \\
\bottomrule
\end{tabularx}
\end{table}

\subsection{算法细节一：回复切分}

模型一次会吐 100 多字，直接发是一大坨。做法是两级切分，目标是「像人打字」，约束是「一个字都不能丢」：

\begin{lstlisting}
// 简化自 src/responder.js：先按行切，再按句末标点切，
// 单句过长时先向后扫句末标点，没有才退到弱标点，最后才硬切
function splitReply(text, max = 100) {
  const out = [];
  for (const line of String(text).split(/\n+/)) {
    let buf = '';
    for (const ch of line) {
      buf += ch;
      if (BUBBLE_END.test(ch)) {           // 。！？…
        if ([...buf].length >= max * 0.5) { out.push(buf.trim()); buf = ''; }
      }
    }
    if (buf.trim()) out.push(buf.trim());
  }
  return out.filter(Boolean);
}
\end{lstlisting}

正确性用「切完拼回去必须等于原文」来断言（测试 \texttt{integration.js}），边界包括：全是标点、超长无标点、UTF-8 代理对（emoji）不能从中间切断。

\subsection{算法细节二：群里该不该说话}

\begin{lstlisting}
// 简化自 src/responder.js：短路判定，每条都是 O(1) 或 O(关键词数)
function shouldReplyGroup(msg, st) {
  if (st.silent) return { skip: 'silent' };                 // 潜水模式
  if (msg.fromOwner && calledMe(msg)) return { reply: true };// 主人在叫我
  if (!st.autoEnabled) return { skip: 'auto-off' };
  if (!hitKeyword(msg.text)) return { skip: 'not-joining' }; // 关键词
  if (Date.now() - st.lastReplyAt < st.minIntervalMs) return { skip: 'cooldown' };
  if (st.repliesThisHour >= st.maxPerHour) return { skip: 'quota' };
  if (Date.now() - msg.timestamp > st.staleSkipMs) return { skip: 'group-stale' };
  return { reply: true };
}
\end{lstlisting}

早期版本「几乎有消息就回」，实测 525 条群消息里它接了 130 次、模型判定不该接的只有 2 次（2.4\%），插话多且答非所问。改成「关键词 + 冷却 + 配额 + 过期」之后，闲聊时 0 发送、被叫到时 1.19 秒内回、过期消息按预期跳过——同一套测试从「不合格」变成「合格」，靠的是把「该不该」写成了可测的判定函数。

\subsection{算法细节三：记忆如何装进有限的上下文}

\begin{itemize}[leftmargin=1.6em]
\item 历史只取最近 \textbf{40 轮}，每条正文截 \textbf{240 字}（转发记录 2000 字、文件 60000 字，走单独通道）；
\item 每 10 轮把对话\textbf{提炼成摘要}，只保留要点，避免历史线性膨胀；
\item 知识库只在「看起来在问事实」时检索，命中才注入前 3 段；
\item 一个容易被忽略的优化：把「当前时间」从系统提示词\textbf{开头}挪到\textbf{末尾}，让前缀保持不变，命中服务端前缀缓存。
\end{itemize}

效果：单次请求提示词从约 \textbf{1082 token 降到约 320 token}，缓存命中率约 \textbf{85\%}。

\subsection{「计算机原理」级别的坑（都进了测试）}

\begin{table}[H]
\centering\footnotesize
\begin{tabularx}{\textwidth}{@{}p{4.3cm}p{5.1cm}X@{}}
\toprule
现象 & 根因 & 处理 \\
\midrule
解压后 \texttt{Cannot find module 'src\textbackslash{}main.js'} & 打包工具用反斜杠做 zip 路径分隔符，Linux 下变成文件名的一部分 & 跨平台统一用 \texttt{/}；自写打包脚本 \\
改完中文 JSON 报 \texttt{Bad control character} & PowerShell 默认按 ANSI 写文件，UTF-8 中文被写坏 & 一律用明确 UTF-8 的 API；读 JSON 先剥 BOM \\
\texttt{.cmd} 报 \texttt{'tart' is not recognized} & 用 LF 写 CRLF 脚本，\texttt{cmd} 行解析错位 & 批处理必须是 CRLF \\
机器人每 5 秒重启一次（909 次/天） & 变量声明前被引用（TDZ），\texttt{ReferenceError} 进崩溃循环 & 加「假 WebSocket 真跑一遍 \texttt{main()}」的启动测试 \\
同一条消息回两遍 & 两个实例都认为自己该回 & 心跳租约锁 + 启动时清理残留 \\
时间全差 8 小时 & 容器时区是 UTC，代码里又叠加了一次偏移 & 统一「内部 UTC / 展示北京时间」 \\
电脑没关机却掉线 & 系统进入现代待机（S0），网卡与长连接同断 & 把电源策略当成系统设计的一部分 \\
\bottomrule
\end{tabularx}
\end{table}

\subsection{性能实测}

\begin{table}[H]
\centering\small
\begin{tabular}{@{}lll@{}}
\toprule
指标 & 数值 & 测量方式 \\
\midrule
私聊首响 & 0.8–1.3\,s & 日志打点（含一次模型调用） \\
群聊平均回复 & 1.07\,s（P90 2.4\,s） & 群聊日志统计 \\
提示词规模 & $\approx$320 token（优化前 1082） & 接口返回的 usage \\
缓存命中率 & $\approx$85\% & usage 的 \texttt{prompt\_cache\_hit} \\
常驻内存 & 稳定，无线性增长 & 环形缓冲 + LRU 上限 \\
测试规模 & 56 文件 / $\approx$1974 断言 & \texttt{node tools/run-tests.js} \\
\bottomrule
\end{tabular}
\end{table}

%=========================================================
\section{测试与结果}
%=========================================================

测试全部是\textbf{可离线运行}的：协议用假 WebSocket，模型用假 fetch，只有 3 个需要真实凭据的用例会打印 SKIP 而不是失败（如实标注，不假装通过）。

\begin{table}[H]
\centering\small
\begin{tabular}{@{}llc@{}}
\toprule
类别 & 代表用例 & 结果 \\
\midrule
协议 & 事件解析、\texttt{msg\_id} 去重、发送失败重试、断线重连 & PASS \\
调度 & 群聊 14 组场景（关键词/冷却/过期/主人/潜水） & PASS \\
记忆 & 40 轮上限、摘要触发、跨天日期、群里非 @ 不记忆 & PASS \\
日程 & 「明天 9 点」「30 分钟后」「周三五点」、取消第 N 条、听不懂给菜单 & PASS \\
检索 & BM25 召回、无关查询不注入 & PASS \\
运维 & 租约锁单实例、发件箱补发、启动冒烟 & PASS \\
边界 & 空消息、纯表情、超长文本、emoji 截断、坏 JSON 配置 & PASS \\
需真实凭据 & 真模型采样 / 私有表情库 & SKIP（3 项） \\
\bottomrule
\end{tabular}
\end{table}

\textbf{可复现}：仓库根目录执行 \texttt{node tools/run-tests.js}，在全新目录 clone 后即可跑（脚本会自动用模板生成占位配置）。实测输出：\texttt{通过 53 跳过 3 失败 0 共 56}。

%=========================================================
\section{AI 协作过程（280 轮）}
%=========================================================

\subsection{整体节奏}

项目不是「让 AI 一次写完」，而是\textbf{需求 → 最小可用 → 真机跑 → 发现新问题 → 再拆}的循环，累计约 280 轮对话。下表是几个代表性阶段：

\begin{table}[H]
\centering\footnotesize
\begin{tabularx}{\textwidth}{@{}p{2.2cm}p{6.2cm}X@{}}
\toprule
阶段 & 我在做什么 & AI 在做什么 \\
\midrule
选题 & 判断官方机器人不可用、逆向不可取，定 OneBot 路线 & 对比方案、列协议细节 \\
最小可用 & 定义「一句话有回复」的验收标准 & 起草接入层、跑通首条链路 \\
拟人化 & 指出「太书面／一次发一大段／群里插话太多」 & 统计真实语料、写切分与调度逻辑 \\
记忆 & 要求「记得住、能查、群里不乱记」 & 设计三层记忆与检索 \\
运维 & 真机掉线、双开、崩溃循环，我盯日志定位 & 读日志、复现、写守护与租约 \\
\bottomrule
\end{tabularx}
\end{table}

\subsection{三个「AI 判断错了、靠验证纠回来」的案例}

\textbf{案例一：双实例不是根因。} 第一次掉线时，AI 从日志里看到两个进程，判断「双实例互踢」并加了单实例守卫。我要求验证：修完仍然掉线，且被踢时刻只有一个实例。于是继续查，最后定位到\textbf{系统进入现代待机}导致长连接断开。结论推翻了上一轮判断——这不是 AI 无能，而是\textbf{必须用实验而不是解释来定案}。

\textbf{案例二：加速反而更慢。} 为降低延迟，曾把「思考链」关掉、把历史压到 6 轮。看上去该更快，实测部分回复质量崩了（答非所问）。改回「群里 6 轮 + 私聊 40 轮」的分层后，延迟与质量才同时可接受。教训是：\textbf{延迟与质量是同一张预算表上的两个科目，不能各改各的}。

\textbf{案例三：真正的 bug 在「看不见的地方」。} 「群里答非所问」我一开始以为是提示词问题，AI 也顺着改提示词。后来我让它统计日志，才发现 84 次准备接话里有 31 次是「排队太久、话题已过」。于是加过期跳过——改的是调度，不是提示词。

\subsection{我承担的工作与验证方法}

\begin{itemize}[leftmargin=1.6em]
\item \textbf{需求与取舍}：方案选择、优先级、什么算「完成」；
\item \textbf{代码审查}：逐段读 AI 生成的代码，改掉不合理的默认值（温度、配额、冷却）；
\item \textbf{真机验证}：部署、盯日志、复现故障、对照实验（改前/改后各跑一次）；
\item \textbf{测试补齐}：把每次踩过的坑写成断言，防止回退。
\end{itemize}

我不把「能跑」当验收标准，而把「能在新环境复现、能在日志里解释」当标准。

%=========================================================
\section{失败案例与已知问题}
%=========================================================

\begin{enumerate}[leftmargin=1.6em]
\item \textbf{掉线三连}：协议端崩溃、双实例互踢、系统现代待机。前两个已修（单实例租约、看护重启），第三个用「保持唤醒 + 电源策略」压住，但笔记本合盖/断电仍可能中断。
\item \textbf{登录态}：整机重启后可能需要重新扫码，这是第三方协议端的固有约束，不是本项目能绕开的。
\item \textbf{模型跑题}：小模型偶尔自说自话，已用低温度 + 输出规范 + 后处理压制，但没有微调，仍有残留。
\item \textbf{知识库能力有限}：只做单跳检索，没有重排与多跳推理。
\item \textbf{无法在评阅环境直接演示 QQ 交互}：需要真实账号与协议端，故本文提供离线测试、日志与截图代替。
\end{enumerate}

%=========================================================
\section{隐私、合规与开源}
%=========================================================

\begin{itemize}[leftmargin=1.6em]
\item \textbf{数据不出本机}：所有记忆与日志落在本地 \texttt{data/}（已在 \texttt{.gitignore}）；只有模型调用出网。
\item \textbf{最小化}：群里非 @ 的消息不写入记忆；日志默认只记长度与耗时。
\item \textbf{可删除}：支持「忘掉 X」「你记住了什么」；删目录即清空。
\item \textbf{凭据}：密钥只在本地 \texttt{config.json}，仓库只有 \texttt{config.example.json} 占位。
\item \textbf{本仓库脱敏}：真实 QQ 号 → 占位数字；昵称 → Alice/小明；聊天记录、表情素材、运维脚本 token 全部移除；测试里涉及私有用例的 3 项改为 SKIP。
\end{itemize}

\textbf{许可证}：MIT。第三方依赖仅 OneBot v11 协议规范、NapCat 协议端、open-meteo、Bing RSS、Noto Emoji（占位素材），均在文档中标注来源。

%=========================================================
\section{总结与展望}
%=========================================================

这个项目让我真正体会到：\textbf{决定一个「AI 应用」好不好用的，往往不是模型，而是模型外面那一圈工程}——去重、限流、过期、缓存、租约、重试、切分、脱敏。课程里的数据结构与算法在这里不是练习题，而是「能不能连续跑三天不崩、会不会把同一条消息回两遍、群聊会不会刷屏」的直接原因。

后续计划：知识库加检索重排；用自建语料微调风格；把运维脚本做成跨平台的一键部署。

\end{document}
